\section{边界：这篇论文没有做什么}
\label{sec:boundary}

前面八章按问题顺序走完了全文。这一章只做一件事：把“本文证明了什么”与“后世补了什么”分开列清。读经典论文最容易犯的错，是把后人补的定理算到作者头上，或者反过来把作者的断言当成已证。

\subsection{本文自己的留白}

\begin{enumerate}[leftmargin=2em]
\item \textbf{一般手术下的标架因子}。§4 对标准手术（模变换 $S$：$S^2\times S^1\to S^3$）算过相位，但对一般整系数 Dehn 手术，“如何在手术下系统地追踪 $e^{2\pi\ii c/24}$ 因子”，原文明确说留待另文 \src{391}。
\item \textbf{函子性}。态空间的对象级等同有 \src{370}，但“沿曲面粘合时共形块如何配对”只有 §5 的 “It is obvious at this point” \src{397}。
\item \textbf{半形式修正的系统处理}。本文只用 $\cL^{\otimes k}$，能级平移以 §2 的 $\eta$ 修正形式出现 \src{359}；以 $K^{1/2}$ 语言给出的独立推导没有。
\item \textbf{“量子群”}。p.372 提到把 $A$ 当经典对象硬解会得到量子群式的东西，但没有构造。
\item \textbf{一般群与一般表示}。文中反复以 $SU(N)$ 定义表示与 $SU(2)$ 为例；对一般单群、一般可积表示，只给框架。
\end{enumerate}

\subsection{后世补了什么}

\begin{center}
\small
\begin{tabular}{@{}p{4.4cm}p{4.2cm}p{4.6cm}@{}}
\toprule
问题 & 本文状态 & 后世 \\
\midrule
三维流形不变量的严格定义 & 手术算法 + 算例；一般手术标架留白 & Reshetikhin--Turaev 1991（组合定义）\\
TQFT 公理化 & 只用切割-粘贴的手工版本 & Atiyah 1988 公理；后世函子化\\
共形块 = 态空间的函子化 & 对象级 + Segal 定理引用 & TUY／后世 TQFT 文献\\
Verlinde 公式的数学证明 & 物理推导（§4.5） & TUY 1989 → Faltings 1994\\
Khovanov 同调 & 不涉及 & Khovanov 2000；物理来源 Witten 2011\\
2+1 维引力 & 姊妹篇 [38] 一句带过 & 与 Brown--Henneaux 结合成 AdS$_3$ 图像（1990s）\\
\bottomrule
\end{tabular}
\end{center}

\subsection{怎样读这篇论文：三遍法}

\begin{enumerate}[leftmargin=2em]
\item \textbf{第一遍（半天）}：读引言、§1、§3 的节首与节末、§5。目标是能说出“作用量是什么、相空间是什么、Hilbert 空间是什么、和共形块什么关系”。
\item \textbf{第二遍（两三天）}：读 §2 与 §4 的推导主干，重点是 (2.16)(2.17) 的 $\eta$ 修正、(4.21) 的 $q$、(4.28)--(4.51) 的手术链。目标是能复算 $S^3$ 上的 Jones 多项式与 $S^2$ 上三点维数。
\item \textbf{第三遍（按需）}：回到具体公式核对约定（$k$ 的归一化、$\Tr$ 的表示、半形式修正），并区分“原文说／我们补／后世”。
\end{enumerate}

\begin{keypoint}
\textbf{一句总结}：这篇论文给的是“一个可精确求解的构造 + 一整套自洽计算 + 若干断言”。把它读成“已经严格的定理集”会处处别扭；把它读成“物理学家给出的、后来被数学严格化的纲领”，每一节的位置就都清楚了。
\end{keypoint}

\begin{boundary}
\begin{itemize}[leftmargin=1.4em]
\item 本章所有“本文状态”均有前文锚点；“后世”一栏的卷页，凡未核者已在正文标出。
\item 未核：RT 1991、TV 1992、Khovanov 2000 的期刊卷页；Atiyah TQFT 公理原文卷页。
\end{itemize}
\end{boundary}
